\chapter{引言}

\section{研究背景}

功能磁共振成像（functional magnetic resonance imaging, fMRI）能够反映大脑在不同任务或疾病状态下的功能活动。与结构影像不同，fMRI 分析关注的不只是单个脑区的局部活动强度，更重要的是不同脑区之间随时间变化而形成的协同活动关系。因此，fMRI 数据常被进一步转化为脑功能连接网络：脑区（region of interest, ROI）作为节点，不同 ROI 之间的功能相关性作为边，进而形成脑图。

在神经影像研究中，如何从高维、噪声较强、样本量相对有限的脑连接数据中提取与疾病或认知状态相关的有效特征，是一个重要问题。传统机器学习方法通常将功能连接矩阵展平成向量后输入分类器，例如 SVM、随机森林或多层感知机。这类方法实现简单，但会削弱脑网络原有的拓扑结构，也难以直接解释哪些脑区或脑区之间的连接对模型判断起到关键作用。

BrainGNN 的提出正是针对这一问题。它将 fMRI 脑功能连接视为图结构数据，并使用图神经网络直接在脑图上进行表示学习。这样既能保留脑区之间的连接关系，又能通过模型内部的节点选择机制为神经影像分析提供一定的可解释性。

\section{现有方法的局限}

早期 fMRI 分类研究多采用传统机器学习方法。这些方法通常需要人工提取特征，或者将连接矩阵展平成高维向量后再进行分类。由于展平操作会破坏脑区之间的网络关系，模型难以直接学习“脑区--连接--整体网络”之间的层级结构。对于样本量有限而特征维度较高的 fMRI 数据，这类方法也容易受到特征冗余和过拟合的影响。

普通深度学习方法虽然能够自动学习复杂非线性特征，但也不一定适合脑功能连接数据。卷积神经网络更适合处理规则网格结构，循环神经网络更适合处理时间序列，而脑功能连接网络本质上是不规则图结构。如果直接使用普通深度学习模型，可能无法充分利用脑区之间的连接拓扑。

传统图论方法长期用于脑网络分析，例如计算节点度、聚类系数、特征路径长度和模块化程度等指标。这些指标具有较强的神经科学解释性，但通常依赖预定义特征，表达能力有限，难以自动学习复杂的多层次脑网络模式。因此，fMRI 脑图分析需要一种同时具备结构建模能力和可解释性的模型。

\section{BrainGNN 的核心思想}

BrainGNN 的核心观点是：脑功能连接数据本质上是一种图结构数据，因此应当使用能够处理图结构的神经网络模型，而不是简单地将连接矩阵展平成向量。BrainGNN 将每个被试的一次 fMRI 数据表示为个体化脑图，并通过图卷积、图池化和读出机制完成图级分类。

与通用图神经网络相比，BrainGNN 更强调脑区的异质性。脑网络中的节点不是匿名节点，不同 ROI 具有固定解剖位置和不同功能意义，因此模型不应当完全共享同一套节点变换规则。BrainGNN 使用 ROI-aware graph convolution（Ra-GConv）为不同 ROI 生成不同的卷积核，使模型能够在学习图结构的同时保留脑区层面的差异。

BrainGNN 的另一个重要设计是 ROI-topK pooling（R-pool）。R-pool 根据节点表示计算每个 ROI 的重要性得分，并保留对分类更有贡献的脑区。该机制一方面可以减少无关或噪声脑区对分类的影响，另一方面也提供了可解释性来源：被模型反复选择或赋予较高权重的 ROI 可以作为潜在 salient ROI 或候选 biomarker 进行分析。

因此，BrainGNN 的贡献不只是提高分类准确率。它更重要的意义在于将脑图结构建模、ROI 特异性和可解释性结合起来，使模型既能完成疾病分类或任务状态识别，也能为理解相关脑区和脑网络模式提供线索。

\section{本报告的工作}

本报告围绕 BrainGNN 的方法、复现、改进和讨论展开。首先，报告总结原论文的实验方法，包括 fMRI 脑图构建、BrainGNN 模型结构、损失函数设计、训练设置和可解释性分析。其次，报告在本地 HCP task-fMRI 数据上进行复现实验，内容包括本地配置说明、与 baseline 方法比较、超参数扫描、卷积层和损失函数消融，以及基于 pooling score 和 community membership 的可解释性图示。

在复现基础上，本报告进一步关注 HCP 数据中的 LR/RL 扫描方向差异。若模型学习到与扫描方向相关而非任务本身相关的差异，可能影响同一受试者、同一任务在不同方向下的预测一致性。因此，报告先后尝试方向对抗训练和 LR/RL 配对一致性训练两种方法，分析它们是否能够在保持任务分类性能的前提下，提升 LR/RL 配对一致性。

最后，报告结合复现和改进实验讨论 BrainGNN 方法贡献、可解释性结果的含义与边界、LR/RL 方向偏移的处理方式，以及本地实验在数据规模、ROI 数量、图构建方式和可解释性验证方面的局限。
